\documentclass[10pt,a4paper]{amsart}
\usepackage{amsmath,amssymb,amsthm,graphicx}
% Retained preamble for the Stacks Project chapter snapshot in this directory.
% The chapter's own preamble.tex is a copy of the shared Stacks Project preamble; the block
% below is sliced MECHANICALLY, line for line, from it. Nothing was added, removed or
% redefined; the wrapper class, the xr-hyper/hyperref wiring and the externaldocument list
% are omitted because the assembler supplies the document class.
%
% lines taken from preamble.tex: 3, 5, 9, 11, 12, 156, 157, 158, ... (51 lines total)
%
% package coverage restored from the chapter's own preamble (2026-10-13): xy, xy-2cell

\usepackage[all]{xy}
\xyoption{2cell}
\UseAllTwocells

\newenvironment{reference}{\comment}{\endcomment}
\newenvironment{slogan}{\comment}{\endcomment}
\newenvironment{history}{\comment}{\endcomment}
\theoremstyle{plain}
\newtheorem{theorem}[subsection]{Theorem}
\newtheorem{proposition}[subsection]{Proposition}
\newtheorem{lemma}[subsection]{Lemma}
\theoremstyle{definition}
\newtheorem{definition}[subsection]{Definition}
\newtheorem{example}[subsection]{Example}
\newtheorem{exercise}[subsection]{Exercise}
\newtheorem{situation}[subsection]{Situation}
\theoremstyle{remark}
\newtheorem{remark}[subsection]{Remark}
\newtheorem{remarks}[subsection]{Remarks}
\numberwithin{equation}{subsection}
\def\lim{\mathop{\mathrm{lim}}\nolimits}
\def\colim{\mathop{\mathrm{colim}}\nolimits}
\def\Spec{\mathop{\mathrm{Spec}}}
\def\Hom{\mathop{\mathrm{Hom}}\nolimits}
\def\Ext{\mathop{\mathrm{Ext}}\nolimits}
\def\SheafHom{\mathop{\mathcal{H}\!\mathit{om}}\nolimits}
\def\SheafExt{\mathop{\mathcal{E}\!\mathit{xt}}\nolimits}
\def\Sch{\mathit{Sch}}
\def\Mor{\mathop{\mathrm{Mor}}\nolimits}
\def\Ob{\mathop{\mathrm{Ob}}\nolimits}
\def\Sh{\mathop{\mathit{Sh}}\nolimits}
\def\NL{\mathop{N\!L}\nolimits}
\def\CH{\mathop{\mathrm{CH}}\nolimits}
\def\proetale{{pro\text{-}\acute{e}tale}}
\def\etale{{\acute{e}tale}}
\def\QCoh{\mathit{QCoh}}
\def\Ker{\mathop{\mathrm{Ker}}}
\def\Im{\mathop{\mathrm{Im}}}
\def\Coker{\mathop{\mathrm{Coker}}}
\def\Coim{\mathop{\mathrm{Coim}}}
\DeclareMathSymbol{\boxtimes}{\mathbin}{AMSa}{"02}
\def\QCohstack{\mathcal{QC}\!\mathit{oh}}
\def\Cohstack{\mathcal{C}\!\mathit{oh}}
\def\Spacesstack{\mathcal{S}\!\mathit{paces}}
\def\Quotfunctor{\mathrm{Quot}}
\def\Hilbfunctor{\mathrm{Hilb}}
\def\Curvesstack{\mathcal{C}\!\mathit{urves}}
\def\Polarizedstack{\mathcal{P}\!\mathit{olarized}}
\def\Complexesstack{\mathcal{C}\!\mathit{omplexes}}
\def\Pic{\mathop{\mathrm{Pic}}\nolimits}
\def\Picardstack{\mathcal{P}\!\mathit{ic}}
\def\Picardfunctor{\mathrm{Pic}}
\def\Deformationcategory{\mathcal{D}\!\mathit{ef}}
\usepackage{cleveref}
\crefname{originallemma}{original lemma}{original lemmas}
\Crefname{originallemma}{Original lemma}{Original lemmas}
\makeatletter
\newlabel{definition-alternative}{{definition-alternative}{source}{}{sourcepointer-1}{}}
\newlabel{definition-alternative@cref}{{[originallemma][0][]definition-alternative}{[1][0][]source}{}{}{}}
\makeatother
\crefname{originallemma}{original lemma}{original lemmas}
\Crefname{originallemma}{Original lemma}{Original lemmas}
\makeatletter
\newlabel{descent-lemma-quasi-coherent-alternative}{{0GZV}{source}{}{sourcepointer-2}{}}
\newlabel{descent-lemma-quasi-coherent-alternative@cref}{{[originallemma][0][]0GZV}{[1][0][]source}{}{}{}}
\makeatother
\crefname{originallemma}{original lemma}{original lemmas}
\Crefname{originallemma}{Original lemma}{Original lemmas}
\makeatletter
\newlabel{equation-alternative-ringed}{{equation-alternative-ringed}{source}{}{sourcepointer-3}{}}
\newlabel{equation-alternative-ringed@cref}{{[originallemma][0][]equation-alternative-ringed}{[1][0][]source}{}{}{}}
\makeatother
\crefname{originallemma}{original lemma}{original lemmas}
\Crefname{originallemma}{Original lemma}{Original lemmas}
\makeatletter
\newlabel{section-restriction}{{section-restriction}{source}{}{sourcepointer-4}{}}
\newlabel{section-restriction@cref}{{[originallemma][0][]section-restriction}{[1][0][]source}{}{}{}}
\makeatother
\crefname{originallemma}{original lemma}{original lemmas}
\Crefname{originallemma}{Original lemma}{Original lemmas}
\makeatletter
\newlabel{sites-modules-definition-site-local}{{03DL}{source}{}{sourcepointer-5}{}}
\newlabel{sites-modules-definition-site-local@cref}{{[originallemma][0][]03DL}{[1][0][]source}{}{}{}}
\makeatother
\crefname{originallemma}{original lemma}{original lemmas}
\Crefname{originallemma}{Original lemma}{Original lemmas}
\makeatletter
\newlabel{sites-modules-lemma-chaotic-quasi-coherent}{{0GZN}{source}{}{sourcepointer-6}{}}
\newlabel{sites-modules-lemma-chaotic-quasi-coherent@cref}{{[originallemma][0][]0GZN}{[1][0][]source}{}{}{}}
\makeatother
\crefname{originallemma}{original lemma}{original lemmas}
\Crefname{originallemma}{Original lemma}{Original lemmas}
\makeatletter
\newlabel{sites-modules-section-intrinsic}{{03DG}{source}{}{sourcepointer-7}{}}
\newlabel{sites-modules-section-intrinsic@cref}{{[originallemma][0][]03DG}{[1][0][]source}{}{}{}}
\makeatother
\title{Stacks, Sheaves on Algebraic Stacks: lemma-quasi-coherent-alternative}
\author{The Stacks Project Authors}
\date{stacks-project:stacks-sheaves.tex:local-snapshot-20261012; GNU FDL 1.2 or later}
\begin{document}
\maketitle
\noindent\textit{Editorial note.} Statement span is lines 4750-4773 of sources/stacks-algebraic-stacks-sheaves/pinned-original-source.tex; proof span is lines 4775-4796, the complete author proof (20 body lines).\par\medskip
\noindent\textit{Editorial note.} Navigation stubs: 7 cited-but-not-included labels are supplied as non-load-bearing pointers with their pinned tags where the snapshot tag index has them.\par\medskip
\noindent\textit{Editorial note.} extra\_packages is empty on purpose: the retained preamble already loads xy with options and a second declaration clashes in the assembler wrapper.\par\medskip
\noindent\textit{Editorial note.} Text is retained verbatim, including the author's own typos and placeholders.\par\medskip
\noindent\textit{Editorial note.} Queue index 20 of queue\_algebraic-stacks-sheaves.json (pb=1260, stmt 4750-4773, proof 4775-4796).\par\medskip
\noindent\textit{Allowed standard prerequisites (editorial).}\begin{itemize}
\item Explicit prerequisite boundary. The cited inputs are prerequisites; the argument in the retained span is not replaced by any cited result.
\end{itemize}\medskip
\ifdefined\hypertarget\hypertarget{sourcepointer-1}{}\fi
\noindent\textit{Source reference.} Original lemma definition-alternative: sources/stacks-algebraic-stacks-sheaves/pinned-original-source.tex, source lines 4639--4646. Cited by the retained statement or proof but not part of the collected spans; supplied only as a navigation pointer. The argument does not depend on its content.\par\medskip
\ifdefined\hypertarget\hypertarget{sourcepointer-2}{}\fi
\noindent\textit{Source reference.} Original lemma 0GZV: sources/stacks-algebraic-stacks-sheaves/pinned-original-source.tex, source lines 4796--4796. Cited by the retained statement or proof but not part of the collected spans; supplied only as a navigation pointer. The argument does not depend on its content.\par\medskip
\ifdefined\hypertarget\hypertarget{sourcepointer-3}{}\fi
\noindent\textit{Source reference.} Original lemma equation-alternative-ringed: sources/stacks-algebraic-stacks-sheaves/pinned-original-source.tex, source lines 4727--4731. Cited by the retained statement or proof but not part of the collected spans; supplied only as a navigation pointer. The argument does not depend on its content.\par\medskip
\ifdefined\hypertarget\hypertarget{sourcepointer-4}{}\fi
\noindent\textit{Source reference.} Original lemma section-restriction: sources/stacks-algebraic-stacks-sheaves/pinned-original-source.tex, source lines 1185--1200. Cited by the retained statement or proof but not part of the collected spans; supplied only as a navigation pointer. The argument does not depend on its content.\par\medskip
\ifdefined\hypertarget\hypertarget{sourcepointer-5}{}\fi
\noindent\textit{Source reference.} Original lemma 03DL: sources/stacks-algebraic-stacks-sheaves/pinned-original-source.tex, source lines 4796--4796. Cited by the retained statement or proof but not part of the collected spans; supplied only as a navigation pointer. The argument does not depend on its content.\par\medskip
\ifdefined\hypertarget\hypertarget{sourcepointer-6}{}\fi
\noindent\textit{Source reference.} Original lemma 0GZN: sources/stacks-algebraic-stacks-sheaves/pinned-original-source.tex, source lines 4796--4796. Cited by the retained statement or proof but not part of the collected spans; supplied only as a navigation pointer. The argument does not depend on its content.\par\medskip
\ifdefined\hypertarget\hypertarget{sourcepointer-7}{}\fi
\noindent\textit{Source reference.} Original lemma 03DG: sources/stacks-algebraic-stacks-sheaves/pinned-original-source.tex, source lines 4796--4796. Cited by the retained statement or proof but not part of the collected spans; supplied only as a navigation pointer. The argument does not depend on its content.\par\medskip
\begin{lemma}
\label{lemma-quasi-coherent-alternative}
Let $p : \mathcal{X} \to (\Sch/S)_{fppf}$ be a category fibred in groupoids.
Let $\mathcal{F}$ be an $\mathcal{O}$-module on $\mathcal{X}_{affine}$.
The following are equivalent
\begin{enumerate}
\item for every morphism $x \to x'$ of $\mathcal{X}_{affine}$ the map
$\mathcal{F}(x') \otimes_{\mathcal{O}(x')} \mathcal{O}(x) \to \mathcal{F}(x)$
is an isomorphism,
\item $\mathcal{F}$ is a quasi-coherent module on
$(\mathcal{X}_{affine}, \mathcal{O})$ in the sense of
Modules on Sites, Definition \ref{sites-modules-definition-site-local},
\item $\mathcal{F}$ is a sheaf for the Zariski topology on
$\mathcal{X}_{affine}$ and a quasi-coherent module on
$(\mathcal{X}_{affine, Zar}, \mathcal{O})$ in the sense of
Modules on Sites, Definition \ref{sites-modules-definition-site-local},
\item same as in (3) for the \'etale topology,
\item same as in (3) for the smooth topology,
\item same as in (3) for the syntomic topology,
\item same as in (3) for the fppf topology, and
\item $\mathcal{F}$ corresponds to a quasi-coherent module
on $\mathcal{X}$ via the equivalence (\ref{equation-alternative-ringed}).
\end{enumerate}
\end{lemma}

\begin{proof}
To make sense out of part (2), recall that $\mathcal{X}_{affine}$ is a site
gotten by endowing the category $\mathcal{X}_{affine}$ with the chaotic
topology (Definition \ref{definition-alternative}) and hence a sheaf
of $\mathcal{O}$-modules $\mathcal{F}$ is the same thing as
a presheaf of $\mathcal{O}$-modules. Conditions (1) and (2)
are equivalent by Modules on Sites, Lemma
\ref{sites-modules-lemma-chaotic-quasi-coherent}.
Observe that for $\tau \in \{Zariski, \etale, smooth, syntomic, fppf\}$
the presheaf $\mathcal{F}$ is a $\tau$-sheaf if and only if
for all $x \in \Ob(\mathcal{X}_{affine})$ the restriction
to $\mathcal{X}_{affine}/x$ is a $\tau$-sheaf. Set $U = p(x)$.
Similarly to the discussion in Section \ref{section-restriction}
the object $x$ of $\mathcal{X}_{affine}$ induces an equivalence
$\mathcal{X}_{affine, \etale}/x \to (\textit{Aff}/U)_\etale$ of sites.
In this way we see that the equivalence of (1) with (3) -- (7)
follows from Descent, Lemma \ref{descent-lemma-quasi-coherent-alternative}
applied to each of these sites. The equivalence of (8) and (7)
is immediate from the fact that ``being quasi-coherent'' is an
intrinsic property of sheaves of modules, see
Modules on Sites, Section \ref{sites-modules-section-intrinsic}
\end{proof}

\end{document}
